% Getaltheorie. Holbein: op de onderste plank ligt een rekenboek voor kooplieden (Peter Apian, 1527),
% opengeslagen bij het delen.
% Motief: het algoritme van Euclides voor 91 en 35 als rechthoek van 91 bij 35, gevuld met zo groot mogelijke
% vierkanten: 91 = 2*35 + 21, 35 = 21 + 14, 21 = 14 + 7, 14 = 2*7. De kleinste vierkanten hebben zijde ggd = 7.
\def\omslagbestand{omslag-ambassadeurs.jpg}
\def\omslagbron{Hans Holbein de Jonge, \textit{De ambassadeurs}, 1533 (publiek domein)}
\def\omslagrand{10.4 + 0.6*\x/21}
\def\omslagmotief{%
  \randlijn
  \begin{scope}[shift={(14.95,7.05)},x=.06cm,y=.06cm]
    \fill[white,fill opacity=.12] (0,0) rectangle (91,35);
    \fill[fel] (84,0) rectangle (91,14);
    \draw[omslaglijn=1pt] (0,0) rectangle (91,35);
    \draw[omslaglijn=1pt] (35,0) -- (35,35) (70,0) -- (70,35) (70,14) -- (91,14) (84,0) -- (84,14) (84,7) -- (91,7);
    \node[text=white,font=\small] at (17.5,17.5) {$35$};
    \node[text=white,font=\small] at (52.5,17.5) {$35$};
    \node[text=white,font=\small] at (80.5,24.5) {$21$};
    \node[text=white,font=\scriptsize] at (77,7) {$14$};
  \end{scope}
  \node[omslaglabel,font=\footnotesize] at (17.68,6.35) {$\mathrm{ggd}(91, 35) = 7$};
}
